\section{Introduction：为何科研、心态与规范}
\label{sec:intro}

\webref{docs/1-导论/}
\resref{resources/01-mindset/}

\subsection{研究动机与问题背景}
做科研与完成课程作业不同：前者面对开放性问题，缺少标准答案与固定教科书路径，主要依赖文献检索、批判性阅读、实验验证与反复迭代。本科/硕士阶段接触科研的价值在于：检验自己是否适应这种节奏、锻炼真实问题解决能力，并为深造或就业积累可迁移资产。与此同时，时间成本高、反馈周期长、方向选择失误都会显著影响体验，因此需要明确预期与规范。

\subsection{能力要求与「三个一定」}
入门阶段尤其需要主动性、自主学习、抗挫折、时间管理与身心健康意识。贯穿全程的三条硬规范是：
\begin{enumerate}
  \item 论文读完必须做笔记（方法、贡献、未解决问题、对自身课题的启发）；
  \item 实验结束后必须留下可回溯记录（环境、配置、数据、结果、分析、下一步）；
  \item 组会必须形成纪要与待办，并将反馈落实到下一阶段任务。
\end{enumerate}

\subsection{常见误区与选题判断}
常见误区包括：盲目迷信顶会顶刊、黑箱调参撞运气、为复杂而复杂、投稿前仍说不清自己解决了什么问题。选题时应回答：问题为何重要、现有方法为何不足、自身资源为何可能解决、能否形成研究主线。有价值的问题往往来自长期接触数据与真实场景，而非简单抄写他人 limitation。换方法可以，应尽量避免频繁更换大方向。论文没有绝对完美版；关键实验稳定、逻辑自洽即可收尾投稿。拒稿是常态，应理性吸收审稿意见。
